% !TEX root = ../../main.tex

\section{Control IO}

El control con \texttt{IO} muestra por qué el reordenamiento no puede activarse para mónadas arbitrarias. En \texttt{IO}, dos acciones que no comparten variables pueden tener efectos observables cuyo orden importa. Por ejemplo, dos impresiones independientes pueden producir salidas distintas si se intercambian.

El log de validación del control negativo registra dos permutaciones. La permutación original imprime:

\begin{verbatim}
A
B
3
\end{verbatim}

Una permutación alternativa imprime:

\begin{verbatim}
B
A
3
\end{verbatim}

La comparación semántica detecta la diferencia y termina con \texttt{RESULT: FAIL}. Este resultado no invalida la solución; al contrario, confirma la necesidad del mecanismo de opt-in. El compilador no debe asumir conmutatividad por ausencia de dependencias RAW, porque los efectos observables pueden seguir dependiendo del orden.

Por esta razón, la extensión solo explora reordenamientos cuando el bloque fue marcado explícitamente como conmutativo mediante \texttt{CD.do}. Para bloques \texttt{do} normales, \texttt{ApplicativeDo} conserva el comportamiento estándar y no se habilita la enumeración de permutaciones conmutativas.
